% !TeX spellcheck = pl_PL
%%%%%%%%%%%%%%%%%%%%%%%%%%%%%%%%%%%%%%%%%%%
%                                        %
% Szablon pracy dyplomowej inzynierskiej %
% zgodny  z aktualnymi  przepisami  SZJK %
%                                        %
%%%%%%%%%%%%%%%%%%%%%%%%%%%%%%%%%%%%%%%%%%
%                                        %
%  (c) Krzysztof Simiński, 2018-2023     %
%                                        %
%%%%%%%%%%%%%%%%%%%%%%%%%%%%%%%%%%%%%%%%%%
%                                        %
% Najnowsza wersja szablonów jest        %
% podstępna pod adresem                  %
% github.com/ksiminski/polsl-aei-theses  %
%                                        %
%%%%%%%%%%%%%%%%%%%%%%%%%%%%%%%%%%%%%%%%%%
%
%
% Projekt LaTeXowy zapewnia odpowiednie formatowanie pracy,
% zgodnie z wymaganiami Systemu zapewniania jakości kształcenia.
% Proszę nie zmieniać ustawień formatowania (np. fontu,
% marginesów, wytłuszczeń, kursywy itd. ).
%
% Projekt można kompilować na kilka sposobów.
%
% 1. kompilacja pdfLaTeX
%
% pdflatex main
% bibtex   main
% pdflatex main
% pdflatex main
%
%
% 2. kompilacja XeLaTeX
%
% Kompilatacja przy użyciu XeLaTeXa różni się tym, że na stronie
% tytułowej używany jest font Calibri. Wymaga to jego uprzedniego
% zainstalowania.
%
% xelatex main
% bibtex  main
% xelatex main
% xelatex main
%
%
%%%%%%%%%%%%%%%%%%%%%%%%%%%%%%%%%%%%%%%%%%%%%%%%%%%%%
% W przypadku pytań, uwag, proszę pisać na adres:   %
%      krzysztof.siminski(małpa)polsl.pl            %
%%%%%%%%%%%%%%%%%%%%%%%%%%%%%%%%%%%%%%%%%%%%%%%%%%%%%
%
% Chcemy ulepszać szablony LaTeXowe prac dyplomowych.
% Wypełniając ankietę spod poniższego adresu pomogą
% Państwo nam to zrobić. Ankieta jest całkowicie
% anonimowa. Dziękujemy!


% https://docs.google.com/forms/d/e/1FAIpQLScyllVxNKzKFHfILDfdbwC-jvT8YL0RSTFs-s27UGw9CKn-fQ/viewform?usp=sf_link
%
%%%%%%%%%%%%%%%%%%%%%%%%%%%%%%%%%%%%%%%%%%%%%%%%%%%%%%%%%%%%%%%%%%%%%%%%%

%%%%%%%%%%%%%%%%%%%%%%%%%%%%%%%%%%%%%%%%%%%%%%%
%                                             %
% PERSONALIZACJA PRACY – DANE PRACY           %
%                                             %
%%%%%%%%%%%%%%%%%%%%%%%%%%%%%%%%%%%%%%%%%%%%%%%

% Proszę wpisać swoje dane w poniższych definicjach.

% TODO
% dane autora
\newcommand{\FirstNameAuthor}{Imię}
\newcommand{\SurnameAuthor}{Nazwisko}
\newcommand{\IdAuthor}{$\langle$wpisać właściwy$\rangle$}   % numer albumu  (bez $\langle$ i $\rangle$)

% drugi autor:
%\newcommand{\FirstNameCoauthor}{Imię}   % Jeżeli jest drugi autor, to tutaj należy podać imię.
%\newcommand{\SurnameCoauthor}{Nazwisko} % Jeżeli jest drugi autor, to tutaj należy podać nazwisko.
%\newcommand{\IdCoauthor}{$\langle$wpisać właściwy$\rangle$}  % numer albumu drugiego autora (bez $\langle$ i $\rangle$)
% Gdy nie ma drugiego autora, należy zostawić poniższe definicje puste, jak poniżej. Gdy jest drugi autor, należy zakomentować te linie.
\newcommand{\FirstNameCoauthor}{} % Jeżeli praca ma tylko jednego autora, to dane drugiego autora zostają puste.
\newcommand{\SurnameCoauthor}{}   % Jeżeli praca ma tylko jednego autora, to dane drugiego autora zostają puste.
\newcommand{\IdCoauthor}{}  % Jeżeli praca ma tylko jednego autora, to dane drugiego autora zostają puste.
%%%%%%%%%%

\newcommand{\Supervisor}{$\langle$tytuł lub stopień naukowy oraz imię i nazwisko$\rangle$}     % dane promotora (bez $\langle$ i $\rangle$)
\newcommand{\Title}{Tytuł pracy dyplomowej inżynierskiej}           % tytuł pracy po polsku
\newcommand{\TitleAlt}{Thesis title in English}                     % thesis title in English
\newcommand{\Program}{$\langle$wpisać właściwy$\rangle$}            % kierunek studiów  (bez $\langle$ i $\rangle$)
\newcommand{\Specialisation}{$\langle$wpisać właściwą$\rangle$}     % specjalność  (bez $\langle$ i $\rangle$)
\newcommand{\Departament}{$\langle$wpisać właściwą$\rangle$}        % katedra promotora  (bez $\langle$ i $\rangle$)

% Jeżeli został wyznaczony promotor pomocniczy lub opiekun, proszę go/ją wpisać ...
\newcommand{\Consultant}{$\langle$stopień naukowy imię i nazwisko$\rangle$} % dane promotora pomocniczego, opiekuna (bez $\langle$ i $\rangle$)
% ... w przeciwnym razie proszę zostawić puste miejsce jak poniżej:
%\newcommand{\Consultant}{} % brak promotowa pomocniczego / opiekuna

% koniec fragmentu do modyfikacji
%%%%%%%%%%%%%%%%%%%%%%%%%%%%%%%%%%%%%%%%%%


%%%%%%%%%%%%%%%%%%%%%%%%%%%%%%%%%%%%%%%%%%%%%%%
%                                             %
% KONIEC PERSONALIZACJI PRACY                 %
%                                             %
%%%%%%%%%%%%%%%%%%%%%%%%%%%%%%%%%%%%%%%%%%%%%%%

%%%%%%%%%%%%%%%%%%%%%%%%%%%%%%%%%%%%%%%%


%%%%%%%%%%%%%%%%%%%%%%%%%%%%%%%%%%%%%%%%%%%%%%%
%                                             %
% PROSZĘ NIE MODYFIKOWAĆ PONIŻSZYCH USTAWIEŃ! %
%                                             %
%%%%%%%%%%%%%%%%%%%%%%%%%%%%%%%%%%%%%%%%%%%%%%%



\documentclass[a4paper,twoside,12pt]{book}
\usepackage[utf8]{inputenc}                                      
\usepackage[T1]{fontenc}  
\usepackage{amsmath,amsfonts,amssymb,amsthm}
\usepackage[british,polish]{babel} 
\usepackage{indentfirst}
\usepackage{xurl}
\usepackage{xstring}
\usepackage{ifthen}
\usepackage{float}



\usepackage{ifxetex}

\ifxetex
	\usepackage{fontspec}
	\defaultfontfeatures{Mapping=tex—text} % to support TeX conventions like ``——-''
	\usepackage{xunicode} % Unicode support for LaTeX character names (accents, European chars, etc)
	\usepackage{xltxtra} % Extra customizations for XeLaTeX
\else
	\usepackage{lmodern}
\fi



\usepackage[margin=2.5cm]{geometry}
\usepackage{graphicx} 
\usepackage{hyperref}
\usepackage{booktabs}
\usepackage{tikz}
\usepackage{pgfplots}
\usepackage{mathtools}
\usepackage{geometry}
\usepackage{subcaption}   % subfigures
\usepackage[page]{appendix} % toc,
\renewcommand{\appendixtocname}{Dodatki}
\renewcommand{\appendixpagename}{Dodatki}
\renewcommand{\appendixname}{Dodatek}

\usepackage{csquotes}
\usepackage[natbib=true,backend=bibtex,maxbibnames=99]{biblatex}  % kompilacja bibliografii BibTeXem
%\usepackage[natbib=true,backend=biber,maxbibnames=99]{biblatex}  % kompilacja bibliografii Biberem
\bibliography{biblio}

\usepackage{ifmtarg}   % empty commands  

\usepackage{setspace}
\onehalfspacing


\frenchspacing



%%%% TODO LIST GENERATOR %%%%%%%%%

\usepackage{color}
\definecolor{brickred}      {cmyk}{0   , 0.89, 0.94, 0.28}

\makeatletter \newcommand \kslistofremarks{\section*{Uwagi} \@starttoc{rks}}
  \newcommand\l@uwagas[2]
    {\par\noindent \textbf{#2:} %\parbox{10cm}
{#1}\par} \makeatother


\newcommand{\ksremark}[1]{%
{%\marginpar{\textdbend}
{\color{brickred}{[#1]}}}%
\addcontentsline{rks}{uwagas}{\protect{#1}}%
}

\newcommand{\comma}{\ksremark{przecinek}}
\newcommand{\nocomma}{\ksremark{bez przecinka}}
\newcommand{\styl}{\ksremark{styl}}
\newcommand{\ortografia}{\ksremark{ortografia}}
\newcommand{\fleksja}{\ksremark{fleksja}}
\newcommand{\pauza}{\ksremark{pauza `--', nie dywiz `-'}}
\newcommand{\kolokwializm}{\ksremark{kolokwializm}}
\newcommand{\cudzyslowy}{\ksremark{,,polskie cudzysłowy''}}

%%%%%%%%%%%%%% END OF TODO LIST GENERATOR %%%%%%%%%%%

\newcommand{\printCoauthor}{%		
    \StrLen{\FirstNameCoauthor}[\FNCoALen]
    \ifthenelse{\FNCoALen > 0}%
    {%
		{\large\bfseries\Coauthor\par}
	
		{\normalsize\bfseries \LeftId: \IdCoauthor\par}
    }%
    {}
} 

%%%%%%%%%%%%%%%%%%%%%
\newcommand{\autor}{%		
    \StrLen{\FirstNameCoauthor}[\FNCoALenXX]
    \ifthenelse{\FNCoALenXX > 0}%
    {\FirstNameAuthor\ \SurnameAuthor, \FirstNameCoauthor\ \SurnameCoauthor}%
	{\FirstNameAuthor\ \SurnameAuthor}%
}
%%%%%%%%%%%%%%%%%%%%%

\StrLen{\FirstNameCoauthor}[\FNCoALen]
\ifthenelse{\FNCoALen > 0}%
{%
\author{\FirstNameAuthor\ \SurnameAuthor, \FirstNameCoauthor\ \SurnameCoauthor}
}%
{%
\author{\FirstNameAuthor\ \SurnameAuthor}
}%

%%%%%%%%%%%% ZYWA PAGINA %%%%%%%%%%%%%%%
% brak kapitalizacji zywej paginy
\usepackage{fancyhdr}
\pagestyle{fancy}
\fancyhf{}
\fancyhead[LO]{\nouppercase{\it\rightmark}}
\fancyhead[RE]{\nouppercase{\it\leftmark}}
\fancyhead[LE,RO]{\it\thepage}


\fancypagestyle{tylkoNumeryStron}{%
   \fancyhf{} 
   \fancyhead[LE,RO]{\it\thepage}
}

\fancypagestyle{bezNumeracji}{%
   \fancyhf{} 
   \fancyhead[LE,RO]{}
}


\fancypagestyle{NumeryStronNazwyRozdzialow}{%
   \fancyhf{} 
   \fancyhead[LE]{\nouppercase{\autor}}
   \fancyhead[RO]{\nouppercase{\leftmark}} 
   \fancyfoot[CE, CO]{\thepage}
}


%%%%%%%%%%%%% OBCE WTRETY  
\newcommand{\obcy}[1]{\emph{#1}}
\newcommand{\english}[1]{{\selectlanguage{british}\obcy{#1}}}
%%%%%%%%%%%%%%%%%%%%%%%%%%%%%

% polskie oznaczenia funkcji matematycznych
\renewcommand{\tan}{\operatorname {tg}}
\renewcommand{\log}{\operatorname {lg}}

% jeszcze jakies drobiazgi

\newcounter{stronyPozaNumeracja}

%%%%%%%%%%%%%%%%%%%%%%%%%%% 
\newcommand{\printOpiekun}[1]{%		

    \StrLen{\Consultant}[\mystringlen]
    \ifthenelse{\mystringlen > 0}%
    {%
       {\large{\bfseries OPIEKUN, PROMOTOR POMOCNICZY}\par}
       
       {\large{\bfseries \Consultant}\par}
    }%
    {}
} 
%
%%%%%%%%%%%%%%%%%%%%%%%%%%%%%%%%%%%%%%%%%%%%%%
 
% Proszę nie modyfikować poniższych definicji!
\newcommand{\Author}{\FirstNameAuthor\ \MakeUppercase{\SurnameAuthor}} 
\newcommand{\Coauthor}{\FirstNameCoauthor\ \MakeUppercase{\SurnameCoauthor}}
\newcommand{\Type}{PROJEKT INŻYNIERSKI}
\newcommand{\Faculty}{Wydział Matematyki Stosowanej} 
\newcommand{\Polsl}{Politechnika Śląska}
\newcommand{\Logo}{politechnika_sl_logo_bw_pion_pl.pdf}
\newcommand{\LeftId}{Nr albumu}
\newcommand{\LeftProgram}{Kierunek}
\newcommand{\LeftSpecialisation}{Specjalność}
\newcommand{\LeftSUPERVISOR}{PROWADZĄCY PRACĘ}
\newcommand{\LeftDEPARTMENT}{KATEDRA}
%%%%%%%%%%%%%%%%%%%%%%%%%%%%%%%%%%%%%%%%%%%%%%

%%%%%%%%%%%%%%%%%%%%%%%%%%%%%%%%%%%%%%%%%%%%%%%
%                                             %
% KONIEC USTAWIEŃ                             %
%                                             %
%%%%%%%%%%%%%%%%%%%%%%%%%%%%%%%%%%%%%%%%%%%%%%%




%%%%%%%%%%%%%%%%%%%%%%%%%%%%%%%%%%%%%%%%%%%%%%%
%                                             %
% MOJE PAKIETY, USTAWIENIA ITD                %
%                                             %
%%%%%%%%%%%%%%%%%%%%%%%%%%%%%%%%%%%%%%%%%%%%%%%

% Tutaj proszę umieszczać swoje pakiety, makra, ustawienia itd.


 
%%%%%%%%%%%%%%%%%%%%%%%%%%%%%%%%%%%%%%%%%%%%%%%%%%%%%%%%%%%%%%%%%%%%%
% listingi i fragmentu kodu źródłowego 
% pakiet: listings lub minted
% % % % % % % % % % % % % % % % % % % % % % % % % % % % % % % % % % % 

% biblioteka listings
\usepackage{listings}
\lstset{%
morekeywords={string,exception,std,vector},% słowa kluczowe rozpoznawane przez pakiet listings
language=C++,% C, Matlab, Python, SQL, TeX, XML, bash, ... – vide https://www.ctan.org/pkg/listings
commentstyle=\textit,%
identifierstyle=\textsf,%
keywordstyle=\sffamily\bfseries, %\texttt, %
%captionpos=b,%
tabsize=3,%
frame=lines,%
numbers=left,%
numberstyle=\tiny,%
numbersep=5pt,%
breaklines=true,%
escapeinside={@*}{*@},%
}

% % % % % % % % % % % % % % % % % % % % % % % % % % % % % % % % % % % 
% pakiet minted
%\usepackage{minted}

% pakiet wymaga specjalnego kompilowania:
% pdflatex -shell-escape main.tex
% xelatex  -shell-escape main.tex

%\usepackage[chapter]{minted} % [section]
%%\usemintedstyle{bw}   % czarno-białe kody 
%
%\setminted % https://ctan.org/pkg/minted
%{
%%fontsize=\normalsize,%\footnotesize,
%%captionpos=b,%
%tabsize=3,%
%frame=lines,%
%framesep=2mm,
%numbers=left,%
%numbersep=5pt,%
%breaklines=true,%
%escapeinside=@@,%
%}

%%%%%%%%%%%%%%%%%%%%%%%%%%%%%%%%%%%%%%%%%%%%%%%%%%%%%%%%%%%%%%%%%%%%%



%%%%%%%%%%%%%%%%%%%%%%%%%%%%%%%%%%%%%%%%%%%%%%%
%                                             %
% KONIEC MOICH USTAWIEŃ                       %
%                                             %
%%%%%%%%%%%%%%%%%%%%%%%%%%%%%%%%%%%%%%%%%%%%%%%



%%%%%%%%%%%%%%%%%%%%%%%%%%%%%%%%%%%%%%%%


\begin{document}
%\kslistofremarks

\frontmatter

%%%%%%%%%%%%%%%%%%%%%%%%%%%%%%%%%%%%%%%%%%%%%%%
%                                             %
% PROSZĘ NIE MODYFIKOWAĆ STRONY TYTUŁOWEJ!    %
%                                             %
%%%%%%%%%%%%%%%%%%%%%%%%%%%%%%%%%%%%%%%%%%%%%%%


%%%%%%%%%%%%%%%%%%  STRONA TYTUŁOWA %%%%%%%%%%%%%%%%%%%
\pagestyle{empty}
{
	\newgeometry{top=1.5cm,%
	             bottom=2.5cm,%
	             left=3cm,
	             right=2.5cm}
 
	\ifxetex 
	  \begingroup
	  \setsansfont{Calibri}
	   
	\fi 
	 \sffamily
	\begin{center}
	\includegraphics[width=50mm]{\Logo}
	 
	
	{\Large\bfseries\Type\par}
	
	\vfill  \vfill  
			 
	{\large\Title\par}
	
	\vfill  
		
	{\large\bfseries\Author\par}
	
	{\normalsize\bfseries \LeftId: \IdAuthor}

	\printCoauthor
	
	\vfill  		
 
	{\large{\bfseries \LeftProgram:} \Program\par} 
	
	{\large{\bfseries \LeftSpecialisation:} \Specialisation\par} 
	 		
	\vfill  \vfill 	\vfill 	\vfill 	\vfill 	\vfill 	\vfill  
	 
	{\large{\bfseries \LeftSUPERVISOR}\par}
	
	{\large{\bfseries \Supervisor}\par}
				
	{\large{\bfseries \LeftDEPARTMENT\ \Departament} \par}
		
	{\large{\bfseries \Faculty}\par}
		
	\vfill  \vfill  

    	
    \printOpiekun{\Consultant}
    
	\vfill  \vfill  
		
    {\large\bfseries  Gliwice \the\year}

   \end{center}	
       \ifxetex 
       	  \endgroup
       \fi
	\restoregeometry
}
  
%%%%%%%%%%%%%%%%%%%%%%%%%%%%%%%%%%%%%%%%%%%%%%%
%                                             %
% KONIEC STRONY TYTUŁOWEJ                     %
%                                             %
%%%%%%%%%%%%%%%%%%%%%%%%%%%%%%%%%%%%%%%%%%%%%%%  


\cleardoublepage

\rmfamily\normalfont
\pagestyle{empty}


%%% No to zaczynamy pisać pracę :-) %%%%

% TODO
\subsubsection*{Tytuł pracy} 
\Title

\subsubsection*{Streszczenie}  
(Streszczenie pracy – odpowiednie pole w systemie APD powinno zawierać kopię tego streszczenia.)

\subsubsection*{Słowa kluczowe} 
(2-5 slow (fraz) kluczowych, oddzielonych przecinkami)

\subsubsection*{Thesis title} 
\begin{otherlanguage}{british}
\TitleAlt
\end{otherlanguage}

\subsubsection*{Abstract} 
\begin{otherlanguage}{british}
(Thesis abstract – to be copied into an appropriate field during an electronic submission – in English.)
\end{otherlanguage}
\subsubsection*{Key words}  
\begin{otherlanguage}{british}
(2-5 keywords, separated by commas)
\end{otherlanguage}




%%%%%%%%%%%%%%%%%% SPIS TRESCI %%%%%%%%%%%%%%%%%%%%%%
% Add \thispagestyle{empty} to the toc file (main.toc), because \pagestyle{empty} doesn't work if the TOC has multiple pages
\addtocontents{toc}{\protect\thispagestyle{empty}}
\tableofcontents

%%%%%%%%%%%%%%%%%%%%%%%%%%%%%%%%%%%%%%%%%%%%%%%%%%%%%
\setcounter{stronyPozaNumeracja}{\value{page}}
\mainmatter
\pagestyle{empty}

\cleardoublepage

\pagestyle{NumeryStronNazwyRozdzialow}

%%%%%%%%%%%%%% wlasciwa tresc pracy %%%%%%%%%%%%%%%%%

% TODO
\chapter{Wstęp}
\label{ch:wstep}

\begin{itemize}
\item wprowadzenie w problem/zagadnienie
\item osadzenie problemu w dziedzinie
\item cel pracy
\item zakres pracy
\item zwięzła charakterystyka rozdziałów
\item jednoznaczne określenie wkładu autora, w przypadku prac wieloosobowych – tabela z autorstwem poszczególnych elementów pracy
\end{itemize}

\section{Wprowadzenie}
Aplikacje społecznościowe w dzisiejszym świecie są dla nas czymś powszechnym. Większość z nich sprowadza się do wysyłania zdjęć na swoje konto aby każdy widział co robisz i gdzie jesteś. Przed erą smartfonów takie informacje były bardziej spersonalizowane, ponieważ służyły do tego listy na których pisaliśmy co robimy, oraz pocztówki dzięki którym mogliśmy pokazać gdzie jesteśmy. Właśnie to drugie zainspirowało nas do połączenia starego z nowym, pocztówek z aplikacją społecznościową. Chcemy aby odwiedzanie różnych miejsc, czy to lokalnie, blisko, czy tych na drugim końcu świata, łączyło się z czymś kiedyś powszechnym czyli wysyłaniem pocztówek. Wszak cieszymy się bardziej z pocztówki wysłanej przez rodzinę lub przyjaciela niż ze zdjęcia na Facebooku.

\section{Osadzenie problemu w dziedzinie}

\section{Cel pracy}
Celem naszej pracy jest stworzenie w pełni działającej aplikacji mobilnej umożliwiającej na zawiązywanie znajomości, wymianę oraz kolekcjonowanie pocztówek. Dodatkowo celem jest stworzenie aplikacji webowej umożliwiającą wykonanie większości czynności z aplikacji mobilnej jak i również backendu działającego na obu platformach odpowiedzialnego za część serwerową.

\section{Zakres pracy}
Zakres pracy obejmuje:
\begin{itemize}
    \item Aplikację mobilną stworzoną w technologi Flutter
    \item Backendu stworzonego w technologi .Net
    \item Aplikcji webowej stworzonej w technologi React
\end{itemize}

\section{Zwięzła charakterystyka rozdziałów}

\section{Autorzy}


% TODO
\chapter{[Analiza tematu]}

\begin{itemize}
\item sformułowanie problemu
\item osadzenie tematu w kontekście aktualnego stanu wiedzy (\english{state of the art}) o poruszanym problemie
\item  studia literaturowe \cite{bib:artykul,bib:ksiazka,bib:konferencja,bib:internet} -  opis znanych rozwiązań (także opisanych naukowo, jeżeli problem jest poruszany w publikacjach naukowych), algorytmów, 
\end{itemize}


Wzory  
\begin{align}
y = \frac{\partial x}{\partial t}
\end{align}
jak i pojedyncze symbole $x$ i $y$  składa się w trybie matematycznym.


%%%%%%%%%%%%%%%%%%%%%%%%





% TODO
\chapter{Wymagania i narzędzia}
\label{ch:wymagania-i-narzedzia}

\begin{itemize}
\item wymagania funkcjonalne i niefunkcjonalne
\item przypadki użycia (diagramy UML) -- dla prac, w których mają zastosowanie
\item opis narzędzi, metod eksperymentalnych, metod modelowania itp.
\item metodyka pracy nad projektowaniem i implementacją -- dla prac, w których ma to zastosowanie
\end{itemize}


% TODO
\chapter{[Właściwy dla kierunku -- np. Specyfikacja zewnętrzna]}
\label{ch:04}

Jeśli „Specyfikacja zewnętrzna”:
\begin{itemize}
\item  wymagania sprzętowe i programowe
\item  sposób instalacji
\item  sposób aktywacji
\item  kategorie użytkowników
\item  sposób obsługi
\item  administracja systemem
\item  kwestie bezpieczeństwa
\item  przykład działania
\item  scenariusze korzystania z systemu (ilustrowane zrzutami z ekranu lub generowanymi dokumentami)
\end{itemize}

%%%%%%%%%%%%%%%%%%%%%
%% RYSUNEK Z PLIKU
%
%\begin{figure}
%\centering
%\includegraphics[width=0.5\textwidth]{./politechnika_sl_logo_bw_pion_pl.pdf}
%\caption{Podpis rysunku zawsze pod rysunkiem.}
%\label{fig:etykieta-rysunku}
%\end{figure}
%Rys. \ref{fig:etykieta-rysunku} przestawia …
%%%%%%%%%%%%%%%%%%%%%
%
%%%%%%%%%%%%%%%%%%%%%
%% WIELE RYSUNKÓW 
%
%\begin{figure}
%\centering
%\begin{subfigure}{0.4\textwidth}
%    \includegraphics[width=\textwidth]{./politechnika_sl_logo_bw_pion_pl.pdf}
%    \caption{Lewy górny rysunek.}
%    \label{fig:lewy-gorny}
%\end{subfigure}
%\hfill
%\begin{subfigure}{0.4\textwidth}
%    \includegraphics[width=\textwidth]{./politechnika_sl_logo_bw_pion_pl.pdf}
%    \caption{Prawy górny rysunek.}
%    \label{fig:prawy-gorny}
%\end{subfigure}
%
%\begin{subfigure}{0.4\textwidth}
%    \includegraphics[width=\textwidth]{./politechnika_sl_logo_bw_pion_pl.pdf}
%    \caption{Lewy dolny rysunek.}
%    \label{fig:lewy-dolny}
%\end{subfigure}
%\hfill
%\begin{subfigure}{0.4\textwidth}
%    \includegraphics[width=\textwidth]{./politechnika_sl_logo_bw_pion_pl.pdf}
%    \caption{Prawy dolny rysunek.}
%    \label{fig:prawy-dolny}
%\end{subfigure}
%        
%\caption{Wspólny podpis kilku rysunków.}
%\label{fig:wiele-rysunkow}
%\end{figure}
%Rys. \ref{fig:wiele-rysunkow} przestawia wiele ważnych informacji, np. rys. \ref{fig:prawy-gorny} jest na prawo u góry.
%%%%%%%%%%%%%%%%%%%%%


 
\begin{figure}
\centering
\begin{tikzpicture}
\begin{axis}[
    y tick label style={
        /pgf/number format/.cd,
            fixed,   % po zakomentowaniu os rzednych jest indeksowana wykladniczo
            fixed zerofill, % 1.0 zamiast 1
            precision=1,
        /tikz/.cd
    },
    x tick label style={
        /pgf/number format/.cd,
            fixed,
            fixed zerofill,
            precision=2,
        /tikz/.cd
    }
]
\addplot [domain=0.0:0.1] {rnd};
\end{axis} 
\end{tikzpicture}
\caption{Podpis rysunku po rysunkiem.}
\label{fig:2}
\end{figure}



% TODO
\chapter{[Właściwy dla kierunku -- np. Specyfikacja wewnętrzna]}
\label{ch:05}

Jeśli „Specyfikacja wewnętrzna”:
\begin{itemize}
\item przedstawienie idei
\item architektura systemu
\item opis struktur danych (i organizacji baz danych)
\item komponenty, moduły, biblioteki, przegląd ważniejszych klas (jeśli występują)
\item przegląd ważniejszych algorytmów (jeśli występują)
\item szczegóły implementacji wybranych fragmentów, zastosowane wzorce projektowe
\item diagramy UML
\end{itemize}

% % % % % % % % % % % % % % % % % % % % % % % % % % % % % % % % % % % 
% Pakiet minted wymaga importu: \usepackage{minted}                 %
% i specjalnego kompilowania:                                       %
% pdflatex -shell-escape main                                       %
% % % % % % % % % % % % % % % % % % % % % % % % % % % % % % % % % % % 


Krótka wstawka kodu w linii tekstu jest możliwa, np.  \lstinline|int a;| (biblioteka \texttt{listings})% lub  \mintinline{C++}|int a;| (biblioteka \texttt{minted})
. 
Dłuższe fragmenty lepiej jest umieszczać jako rysunek, np. kod na rys \ref{fig:pseudokod:listings}% i rys. \ref{fig:pseudokod:minted}
, a naprawdę długie fragmenty – w załączniku.


\begin{figure}[H]
\begin{lstlisting}
class test : public basic
{
    public:
      test (int a);
      friend std::ostream operator<<(std::ostream & s, 
                                     const test & t);
    protected:
      int _a;  
      
};
\end{lstlisting}
\caption{Pseudokod w \texttt{listings}.}
\label{fig:pseudokod:listings}
\end{figure}

%\begin{figure}
%\centering
%\begin{minted}[linenos,frame=lines]{c++}
%class test : public basic
%{
%    public:
%      test (int a);
%      friend std::ostream operator<<(std::ostream & s, 
%                                     const test & t);
%    protected:
%      int _a;  
%      
%};
%\end{minted}
%\caption{Pseudokod w \texttt{minted}.}
%\label{fig:pseudokod:minted}
%\end{figure}




\section{Biblioteki}
\subsection{Biblioteka główna mobilna}
\begin{itemize}
    \item Flutter - jest to technologia stworzona przez firmę Google pozwalająca na tworzenie aplikacji na urządzenia Android jak i iOS używając tego samego kodu. Językiem wykorzystywanym podczas tworzenia tych aplikacji jest Dart który również został stworzony przez Google. Napisane we Flutterze aplikacje mobilne wyglądają i działają niemal identycznie na obu systemach z wydajnością która jest bardzo zbliżona do aplikacji natywnych. Fluttera tworzą dwa podstawowe elementy którymi są framework który jest niezbędny do poprawnego działania aplikacji oraz pakiet SDK który jest konieczny do ich tworzenia.
\end{itemize}
\subsection{Biblioteki mobilne dodatkowe}
\begin{itemize}
    \item intl - umożliwia stworzenie pliku zawierającego wszystkie wiadomości zawierające się w aplikacji oraz używanie odpowiednich plików z językami w zależności od tego wybranego w ustawieniach.
    \item http - jest to biblioteka umożliwiająca komunikację pomiędzy aplikacją a serwerem używając żądań HTTP.
    \item flutter\_bloc - biblioteka która zawiera w sobie widgey ułatwiająca wprowadzneie BLoC jako state management w naszej aplikacji.
    \item flutter\_secure\_storage - w aplikacji używany jest do przechowywania JWT otrzymanego podczas logowania. Biblioteka szyfruje wartości jakie otrzymuje przed umieszczeniem ich w pamięci. Na platformie Android jest używane do tego Szyfrowanie AES, natomiast na platformie iOS jest to KeyChain. 
    \item jwt\_decoder - jest małą biblioteka do dekodowania JWT. Dzięki niej możemy sprawdzić, czy token wygasł i na podstawie tego przekierować użytkownika na odpowiednią podstronę.
    \item shimmer - biblioteka dodająca widget używany podczas ładowania danych z serwera.
    \item image\_picker - oficjalna biblioteka umożliwiająca wybieranie zdjęć z galerii lub tworzenie od razu nowych i dalsze ich wykorzystywanie w aplikacji.
    \item country\_picker - biblioteka dodająca widget do wybierania kraju.
    \item auto\_size\_text - jest to biblioteka dająca większą kontrolę nad wyświetlanym tekstem. Dzięki niej łatwo możemy dla przykładu skalować tekst na mniejszych ekranach lub ustawiać limit wierszy przy zawijaniu.
    \item background\_locator\_2 - jest to jedna z istotniejszych bibliotek w naszej aplikacji. Służy ona na do pobierania z telefonu koordynatów. Działanie funkcji nie jest przerywane przez danie aplikacji do karty lub jej całkowite wyłączenie. Biblioteka jest również rozwinięta o funkcje callback która jest wykonywana za każdym razem gdy pobierany są nowe koordynaty co umożliwia nam wykonywanie wtedy innych funkcji, nawet gdy aplikacja jest wyłączona.
    \item path\_provider - biblioteka wykorzystywana na lokalizowanie plików w systemie.
    \item location\_permissions - umożliwia pobieranie lokalizacji z urządzenia. 
    \item permission\_handler- biblioteka która obsługuje nadawanie oraz sprawdzanie istniejących zgód na wykorzystywanie poszczególnych serwisów w urządzeniu takich jak powiadomienia lub pamięć. 
    \item geolocator - biblioteka która pozwala na komunikacje z serwisem odpowiedzialnym za lokalizację. W projekcie jest wykorzystywana w celu sprawdzenia czy lokazacja jest ciągle włączona na urządzeniu. 
    \item cached\_memory\_image - biblioteka umożliwiająca na zapisywanie grafik w pamięci telefonu co eliminuje miganie obrazów podczas ich ciągłego pobierania.
    \item url\_launcher - umożliwia na otwieranie linków w osobnej karcie przeglądarki na użądzeniach mobilnych. 
    \item loading\_animation\_widget - biblioteka zawierająca animacje ładowania.
    \item flutter\_animate - jest to biblioteka umożliwiająca animowania dowolnego widgetu w naszej aplikacji.
    \item easy\_debounce - biblioteka uproszczająca oryginalny zapis debouncera języka Dart. 
    \item shared\_preferences - służy do zapisywania różnych zmiennych w pamięci telefonu jako pary klucz-wartość w podobny sposób jak to robi secure storage. 
    \item app\_settings - umożliwia otworzenie ustawień telefonu z poziomu aplikacji
\end{itemize}
\section{Architektura systemu}
\subsection{Architektura modułu mobilnego}
Architektura oparta na zdarzeniach (Event-Driven Architecture, EDA) stanowi kluczowy element naszego projektu aplikacji mobilnej. Wybraliśmy tę architekturę ze względu na jej zdolność do efektywnego zarządzania komunikacją między różnymi modułami i komponentami aplikacji, zwłaszcza w kontekście obsługi stanu i przetwarzania zdarzeń.

W ramach tej architektury, BLoC i Cubit pełnią rolę obsługiwaczy zdarzeń.
BLoC to jedno z narzędzi do zarządzania stanem w aplikacjach Flutter. Możemy go używać do obsługi wszystkich stanów, które chcemy kontrolować w naszych aplikacjach. Cubit to podzbiór wzorca projektowego BLoC, który upraszcza sposób zarządzania stanem aplikacji. W tym celu zastępuje użycie zdarzeń (używanych w BLoC) funkcjami, które odbudowują interfejs użytkownika poprzez emitowanie różnych stanów na strumieniu.

Cubit jest podobny do BLoC, ale nie posiada pojęcia o zdarzeniach i polega na metodach do emitowania nowych stanów. Są odpowiedzialne za nasłuchiwanie i przetwarzanie zdarzeń, które są generowane w różnych częściach aplikacji. Na przykład, BLoC lub Cubit ds. uwierzytelniania może nasłuchiwać zdarzeń takich jak "zalogowanie użytkownika" lub "wylogowanie użytkownika" i reagować poprzez zmianę stanu uwierzytelnienia.

Zdarzenia są generowane w odpowiedzi na różne interakcje użytkownika oraz na operacje asynchroniczne. Przykłady to zdarzenia "kliknięcia przycisku" w interfejsie użytkownika, "wysłania formularza" lub "zakończenia operacji pobierania danych z serwera". Zdarzenia te są kluczowe dla przekazywania informacji o działaniach i stanach aplikacji.

Jednym z głównych celów EDA jest utrzymanie luźnego powiązania między różnymi komponentami naszej aplikacji. Oznacza to, że BLoC i Cubit nie muszą mieć wiedzy o źródle generowania zdarzenia, co przyczynia się do modularności i skalowalności projektu. Nowe funkcje i modyfikacje istniejących funkcji mogą być wprowadzane bez konieczności naruszania istniejącej struktury.

Aplikacje Flutter często wymagają operacji asynchronicznych, takich jak komunikacja z serwerem lub dostęp do bazy danych. Dzięki architekturze opartej na zdarzeniach, te operacje asynchroniczne są inicjowane i obsługiwane poprzez generowanie i przetwarzanie zdarzeń. Jest to ważne dla zachowania spójności i zarządzania asynchronicznymi przepływami w aplikacji.

Architektura oparta na zdarzeniach umożliwia także wprowadzenie warstwy pośredniej, która może przechwytywać, przetwarzać i filtrować zdarzenia przed ich dostarczeniem do odpowiednich BLoC lub Cubit. To zapewnia dodatkową kontrolę nad przepływem zdarzeń i może być używane do wykonywania ogólnych operacji, takich jak logowanie lub uwierzytelnianie.

Wszystkie te elementy razem tworzą spójną i skalowalną architekturę, która umożliwia nam efektywne zarządzanie stanem i komunikacją w naszej aplikacji.

\subsection{Struktura plików}
\subsubsection{Aplikacja mobilna}

Poniżej znajdują się nazwy głównych folderów oraz opis plików jakie w sobie posiadają:
\begin{itemize}
    \item api - znajdują się tam wszystkie modele związane z tworzeniem zapytań oraz odbieraniem danych z serwera. Każdy model request posiada metodę toJson która zamienia obiekt flutterowy na JSON'a, natomiast modele respons, posiadają metodę fromJson która zamienia JSON'a który został nam przesłany z serwera na obiekt flutterowy który możemy wykorzystać w naszej aplikacji.

    \item constants - folder ten zawiera wszelkie stałe które są używane w wielu miejscach aplikacji dla przykładu baseUrl który jest niezmienną częścią adresu url do połączenia z backendem.

    \item cubit - zawiera on pliki dotyczące zarządzania stanem aplikacji. Dzięki użyciu cubita zamiast czystego BLoCa nie posiadamy w nim plików związanych z eventami ponieważ funkcjonalność eventów przejęły funkcje znajdujące się już w cubitach. Poza cubitami które zmieniają stany aplikacji oraz obsługują logikę, folder ten posiada również pliki stanów które określają jakie zmienne posiadają konkretne stany aplikacji.

    \item custom\_widgets - jest to zbiór naszych ostylowanych już widgetów przygotowanych do wielokrotnego użytku. W większości sprowadza się to do tego, że dany widget ma własności które posiadają wartości domyślne które możemy zmienić w zależności od potrzeb lub zostawić takimi jakimi są, lecz są również własności które należy podać.

    \item extensions - są to pliki które umożliwiają rozwijanie wbudowanych klas takich jak na przykład String o nowe metody które nie są domyślnie zawarte we frameworku takie jak na przykład Ustawianie wielkiej litery na początku.

    \item helpers - są tu zawarte funkcje które z uwagi na wielokrotne użycie w wielu miejscach naszej aplikacji, zostały przeniesione do osobnego folderu. Znajdują się tu dla przykładu funkcje wyświetlające snackbary lub formatujące datę zgodnie z ustawieniami użytkownika. 

    \item l10n - jest to folder zawierający pliki .arb które są odpowiedzialne za przechowywanie tekstów znajdujących się w aplikacji. Dzięki zawarciu wszystkich tekstów tutaj, tłumaczenie aplikacji sprowadza się do dodania nowego pliku .arb zawierającego wszystkie teksty, oraz dodanie obsługiwanego tekstu do main.dart. 

    \item pages - są tutaj wszystkie strony jakie posiadamy w aplikacji. są one przechowywane w osobnych folderach ponieważ każda ze stron jest rozdzielona na mniejsze widgety które z uwagi na swoje jednokrotne zastosowanie i brak uniwersalności nie możemy przechowywać w folderze custom\_widgets. 

    \item services - znajdują się tutaj pliki odpowiedzialne za komunikację pomiędzy aplikacją a serwerem, jak i również komunikację wewnątrz aplikacji jak na przykład pobieranie danych GPS urządzenia czy dostęp do pamięci wewnętrznej urządzenia. 
\end{itemize}

\subsection{Komunikacja aplikacji mobilnej z serwerem}
Komunikacja aplikacji mobilnej z serwerem jest przeprowadzana za pomocą serwisów oraz cubitów. Cubity obsługują eventy które są wywoływane przez użytkownika. Obsługa eventów w większości przypadków przebiega następująco:
\begin{itemize}
    \item Zamiana stanu na stan ładownia.
    \item Jeżeli jest to event polegający na wysłaniu zapytania do serwera to teraz jest ono wykonywane.
    \item Przetworzenie informacji zwrotnych z serwera.
    \item Ustawienie stanu załadowanego zawierającego nowe dane.
\end{itemize}

Wysłanie zapytania tak jak to było wspomniane powyżej jest wykonywane przez serwis. Każdy serwis zawiera zestaw zapytań dotyczący jednej dziedziny takiej jak autoryzacja lub dane szczegółowe użytkownika. Każde zapytanie poza rejestracją oraz logowaniem wymaga posiadania tokenu. Token ten otrzymujemy podczas logowania i przechowujemy go w secured storage. Jest on pobierany z pamięci urządzenia za każdym razem gdy wykonujemy zapytanie. Odpowiedź z serwera jest sprawdzana i w zależności czy zapytanie powiodło się to wywoływana jest metoda fromJson z odpowiedniego modelu response lub wyświetlany jest komunikat w postaci snackbara.

% TODO
\chapter{Weryfikacja i walidacja}
\label{ch:06}
\begin{itemize}
\item sposób testowania w ramach pracy (np. odniesienie do modelu V)
\item organizacja eksperymentów
\item przypadki testowe zakres testowania (pełny/niepełny)
\item wykryte i usunięte błędy
\item opcjonalnie wyniki badań eksperymentalnych
\end{itemize}

\begin{table}
\centering
\caption{Nagłówek tabeli jest nad tabelą.}
\label{id:tab:wyniki}
\begin{tabular}{rrrrrrrr}
\toprule
	         &                                     \multicolumn{7}{c}{metoda}                                      \\
	         \cmidrule{2-8}
	         &         &         &        \multicolumn{3}{c}{alg. 3}        & \multicolumn{2}{c}{alg. 4, $\gamma = 2$} \\
	         \cmidrule(r){4-6}\cmidrule(r){7-8}
	$\zeta$ &     alg. 1 &   alg. 2 & $\alpha= 1.5$ & $\alpha= 2$ & $\alpha= 3$ &   $\beta = 0.1$  &   $\beta = -0.1$ \\
\midrule
	       0 &  8.3250 & 1.45305 &       7.5791 &    14.8517 &    20.0028 & 1.16396 &                       1.1365 \\
	       5 &  0.6111 & 2.27126 &       6.9952 &    13.8560 &    18.6064 & 1.18659 &                       1.1630 \\
	      10 & 11.6126 & 2.69218 &       6.2520 &    12.5202 &    16.8278 & 1.23180 &                       1.2045 \\
	      15 &  0.5665 & 2.95046 &       5.7753 &    11.4588 &    15.4837 & 1.25131 &                       1.2614 \\
	      20 & 15.8728 & 3.07225 &       5.3071 &    10.3935 &    13.8738 & 1.25307 &                       1.2217 \\
	      25 &  0.9791 & 3.19034 &       5.4575 &     9.9533 &    13.0721 & 1.27104 &                       1.2640 \\
	      30 &  2.0228 & 3.27474 &       5.7461 &     9.7164 &    12.2637 & 1.33404 &                       1.3209 \\
	      35 & 13.4210 & 3.36086 &       6.6735 &    10.0442 &    12.0270 & 1.35385 &                       1.3059 \\
	      40 & 13.2226 & 3.36420 &       7.7248 &    10.4495 &    12.0379 & 1.34919 &                       1.2768 \\
	      45 & 12.8445 & 3.47436 &       8.5539 &    10.8552 &    12.2773 & 1.42303 &                       1.4362 \\
	      50 & 12.9245 & 3.58228 &       9.2702 &    11.2183 &    12.3990 & 1.40922 &                       1.3724 \\
\bottomrule
\end{tabular}
\end{table}  



% TODO
\chapter{Podsumowanie i wnioski}
\begin{itemize}
\item uzyskane wyniki w świetle postawionych celów i zdefiniowanych wyżej wymagań
\item kierunki ewentualnych danych prac (rozbudowa funkcjonalna …)
\item problemy napotkane w trakcie pracy
\end{itemize}



\backmatter

%\bibliographystyle{plplain}  % bibtex
%\bibliography{biblio} % bibtex
\printbibliography           % biblatex
\addcontentsline{toc}{chapter}{Bibliografia}

\begin{appendices}

% TODO
\chapter{Spis skrótów i symboli}

\begin{itemize}
\item[DNA] kwas deoksyrybonukleinowy (ang. \english{deoxyribonucleic acid})
\item[MVC] model -- widok -- kontroler (ang. \english{model--view--controller}) 
\item[$N$] liczebność zbioru danych
\item[$\mu$] stopnień przyleżności do zbioru
\item[$\mathbb{E}$] zbiór krawędzi grafu
\item[$\mathcal{L}$] transformata Laplace'a 
\end{itemize}


% TODO
\chapter{Źródła}

Jeżeli w pracy konieczne jest umieszczenie długich fragmentów kodu źródłowego, należy je przenieść w to miejsce.

\begin{lstlisting}
if (_nClusters < 1)
	throw std::string ("unknown number of clusters");
if (_nIterations < 1 and _epsilon < 0)
	throw std::string ("You should set a maximal number of iteration or minimal difference -- epsilon.");
if (_nIterations > 0 and _epsilon > 0)
	throw std::string ("Both number of iterations and minimal epsilon set -- you should set either number of iterations or minimal epsilon.");
\end{lstlisting}


% % % % % % % % % % % % % % % % % % % % % % % % % % % % % % % % % % % 
% Pakiet minted wymaga odkomentowania w pliku config/settings.tex   %
% importu pakietu minted: \usepackage{minted}                       %
% i specjalnego kompilowania:                                       %
% pdflatex -shell-escape praca                                      %
% % % % % % % % % % % % % % % % % % % % % % % % % % % % % % % % % % % 

%\begin{minted}[linenos,breaklines,frame=lines]{c++}
%if (_nClusters < 1)
%   throw std::string ("unknown number of clusters");
%if (_nIterations < 1 and _epsilon < 0)
%   throw std::string ("You should set a maximal number of iteration or minimal difference -- epsilon.");
%if (_nIterations > 0 and _epsilon > 0)
%   throw std::string ("Both number of iterations and minimal epsilon set -- you should set either number of iterations or minimal epsilon.");
%\end{minted}


% TODO
\chapter{Lista dodatkowych plików, uzupełniających tekst pracy} 


W systemie do pracy dołączono dodatkowe pliki zawierające:
\begin{itemize}
\item źródła programu,
\item dane testowe,
\item film pokazujący działanie opracowanego oprogramowania lub zaprojektowanego i~wykonanego urządzenia,
\item itp.
\end{itemize}


\listoffigures
\addcontentsline{toc}{chapter}{Spis rysunków}
\listoftables
\addcontentsline{toc}{chapter}{Spis tabel}

\end{appendices}

\end{document}


%% Finis coronat opus.

